% ---------------------------------------------------------------------
%  4 UŽDUOTIS — DI pagrįsto aptikimo sprendimo kūrimas
%  REIKALINGA PREAMBULĖJE: xltabular, tabularx, booktabs, siunitx,
%                          graphicx, float
%  Lentelės ir paveikslai GENERUOJAMI:
%    python -m src.eksperimentai.lenteles     -> lenteles/imtis.tex
%    python -m src.eksperimentai.i_latex      -> lenteles/veikimas.tex
%    python -m src.eksperimentai.slenkstis    -> lenteles/slenkstis.tex
%    python -m src.eksperimentai.paveikslai   -> paveikslai/architektura.pdf
% ---------------------------------------------------------------------

\subsection{Sprendimo architektūra}
\label{sec:architektura}

Sprendimas kuriamas kraštiniam šliuzui.

Aptikimo grandinė ir jos paruošimas atskirti (\ref{pav:architektura}
paveikslas). Šliuze srauto langas, 10 arba 100 paketų agregatas,
paverčiamas 36 požymių vektoriumi; neuroniniams modeliams jis dar
normalizuojamas; modelis grąžina atakos įvertį; sprendimas priimamas
palyginus įvertį su slenksčiu. Pats modelis ir slenkstis paruošiami ne
šliuze.

\begin{figure}[H]
  \centering
  \includegraphics[width=\textwidth]{architektura.pdf}
  \caption{Aptikimo sprendimo architektūra}
  \label{pav:architektura}
\end{figure}

Slenkstis nustatomas iš validacijos aibės, ne iš diegimo srauto. Pagal
stebimą srautą kalibruojama sistema prisitaikytų prie atakos, jei ši
truktų pakankamai ilgai.


\subsection{Duomenų paruošimas}
\label{sec:duomenu_paruosimas}

Iš \num{45019243} rinkinio eilučių pašalinamos nutrūkusios ir turinčios
begalines \texttt{Rate} reikšmes, iš viso \num{1000}. Toliau šalinami
tikslūs dublikatai, ir tik po to sudaroma imtis.

\begin{table}[htbp]
\centering
\caption{Duomenų rinkinys pagal kategorijas: dublikatų dalis, unikalių
         eilučių skaičius ir imtis}
\label{tab:imtis}
\lentele{imtis.tex}
\end{table}

Dublikatai pasiskirstę netolygiai (\ref{tab:imtis} lentelė): potvynio
klasėse jų yra dešimtys procentų, o gerybiniame sraute ir retose atakų
klasėse beveik nėra. Todėl dublikatų šalinimas pats savaime mažina
disbalansą, nes traukiasi gausiausios klasės.

Imtis sudaroma taikant \num{100000} eilučių ribą kiekvienai etiketei;
retos klasės imamos visos. Rezultatas yra \num{2425937} eilutės, arba
5,4~\% valyto rinkinio. Trylikai klasių iš 34 riba neįsijungia, o
mažiausia klasė \texttt{UPLOADING\_ATTACK} turi \num{1196} eilutes.

Kadangi riba taikoma etiketei, kategorijų dydžiai imtyje lieka labai
nevienodi. Disbalansas imtyje yra 84:1; be ribos jis būtų 5764:1.

Prieštaringos etiketės, t.\,y. vienodi požymių vektoriai, pažymėti
skirtingai, paliekamos kaip tikras duomenų dviprasmiškumas. Imtyje jų yra
\num{5114}, arba \num{10435} eilutės (0,43~\%); teorinė tikslumo riba
šioje imtyje yra 99,78~\%.

Skaidymas yra stratifikuotas atsitiktinis 70/15/15: \num{1698155}
eilutės mokymui, po \num{363891} validacijai ir testavimui.
Stratifikuojama pagal 34 etiketes, nes kategorija yra etiketės funkcija,
tad kategorijų proporcijos išlieka savaime. Didžiausias klasės
proporcijos nuokrypis mokymo aibėje siekia 0,0001 procentinio punkto.


\subsection{Požymių atranka}
\label{sec:pozymiai}

Iš 39 leidimo požymių šalinami trys, kurių ryšys su kitais yra tikslus:
\texttt{Variance} (= \texttt{Std}$^2$), \texttt{Tot sum}
(= \texttt{AVG} $\times$ \texttt{Number}) ir \texttt{Tot size}
(= \texttt{AVG}). Visos trys tapatybės
patikrintos darbinėje imtyje, nesutapimų nė vieno iš \num{2425937}.
Lieka 36 požymiai.

Šalinama sąrašu, ne automatiniu filtru. Koreliacijos slenkstis šių porų
nepagautų, nes \texttt{Variance} ir \texttt{Std} Pearson koreliacija
tėra \num{0.737}, nors ryšys tikslus. Mažos dispersijos filtras
pašalintų šešis požymius, kurie daugiau nei 99,5~\% eilučių turi tą
pačią reikšmę, bet rečiausiose klasėse įgyja 7--22 kartus didesnes
vidutines reikšmes, o tos klasės ir lemia macro-F1.

\texttt{Protocol Type} paliekamas, nors atrodo perteklinis šalia
\texttt{TCP}, \texttt{UDP} ir \texttt{ICMP} stulpelių. Sutapimas su
jais yra 69,8--94,4~\%, t.\,y. koreliacija, o ne tapatybė.

Vėliavėlių stulpeliai (\texttt{syn\_flag\_number} ir kiti) nėra
dvejetainiai, tai vėliavėlių dalys lange.


\subsection{Modelių realizacija}
\label{sec:sasaja}

Visi keturi metodai realizuoti su vienoda sąsaja: \texttt{fit},
\texttt{predict}, \texttt{predict\_proba} ir išsaugojimas. Du sąsajos
laukai nurodo, ar modelis prižiūrimas ir ar jam būtinas normalizavimas.

\texttt{predict\_proba} grąžina įvertį, kurio didesnė reikšmė reiškia
didesnę atakos tikimybę. Prižiūrimiems metodams tai tikimybių matrica,
autokoderiui atkūrimo paklaida; bendras yra tik rikiavimas, tad
skaičiuojami PR-AUC ir ROC-AUC.

Normalizavimui naudojamas \texttt{StandardScaler}, pritaikomas tik
mokymo aibei. Skalė išsaugoma kartu su modeliu.

Klasių disbalansas sprendžiamas klasių svoriais; jų santykis mokymo
aibėje yra 83,9. XGBoost atveju naudojamas eilučių svorių masyvas, nes
\texttt{scale\_pos\_weight} daugiaklasėje užduotyje neveikia. Gerybinis
srautas nesintetinamas, nes klaidingų teigiamų analizė turi remtis tikru
srautu.


\subsection{Sprendimo slenkstis}
\label{sec:slenkstis}

Prižiūrimi modeliai grąžina tikimybes, tad sprendimo taškas parenkamas
po mokymo. Prie numatytojo pasirinkimo, kai atsakymu imama labiausiai tikėtina
klasė (lentelėse žymima \texttt{argmax}), klaidingų teigiamų dalis
siekia 21--31~\%, t.\,y. 21--31 kartus viršija
\ref{sec:atakos} skyriuje priimtą biudžetą.

Todėl ataka skelbiama tik tada, kai bendra atakų tikimybė viršija
slenkstį $\tau$. Slenkstis parenkamas kaip mažiausias $\tau$,
tenkinantis klaidingų teigiamų biudžetą, ir renkamas tik validacijos
aibėje; išmatuoti operaciniai taškai pateikiami
\ref{sec:slenkscio_kompromisas} poskyryje.

Biudžeto kaina modeliams nevienoda, ir priežastis yra tikimybių
skiriamoji geba. Random Forest tikimybė yra medžių balsų dalis, todėl
ties vienetu ji grubi, o stiprinimo metodas grąžina tolydžias reikšmes
ir aukštą slenkstį pasiekia tiksliau.

Autokoderiui ši taisyklė netaikoma, nes jis klasių neturi. Jo
operacinis taškas yra atkūrimo paklaidos procentilis iš validacijos
aibės gerybinio srauto. Pasirinktas 99-asis procentilis atitinka tą patį 1~\% biudžetą.


\subsection{Eksperimentų infrastruktūra ir prototipas}
\label{sec:infrastruktura}

Kiekvienas eksperimentas aprašomas konfigūracijos failu, kuris nurodo
modelį, užduoties formuluotę ir hiperparametrus; rezultatas yra viena
eilutė bendrame rezultatų faile, iš kurio generuojamos ataskaitos
lentelės.

Du nepriklausomi pilni perleidimai davė kokybės metrikas, sutampančias
iki $1\cdot10^{-5}$.

\begin{table}[htbp]
\centering
\caption{Veikimo rodikliai: mokymo laikas, inferencijos delsa procesoriumi
         ir modelio dydis}
\label{tab:veikimas}
\lentele{veikimas.tex}
\end{table}

Inferencijos delsa (\ref{tab:veikimas} lentelė) matuojama procesoriumi,
nors mokymui naudotas grafinis spartintuvas, nes kraštinis šliuzas jo
neturi.

Sprendimo veikimui pademonstruoti sukurtas prototipas
(\ref{pav:prototipas} paveikslas): rodomi apdorotų langų ir signalų
skaičiai, klaidingų teigiamų dalis, aptikimo dalis pagal kategorijas ir
inferencijos delsa. Prototipas naudoja
tą patį sprendimo tašką kaip vertinimas ir dirba su validacijos aibe.

\begin{figure}[H]
  \centering
  \includegraphics[width=\textwidth]{prototipas.png}
  \caption{Prototipas apdorojant \num{5000} srauto langų. Klaidingų
    teigiamų dalis yra 1,29~\%, aptinkama 89,6~\% atakų. Bendras rodiklis
    slepia skirtumus tarp kategorijų. Potvynio ir Mirai atakos aptinkamos
    visos, o \texttt{BruteForce} ir žiniatinklio atakos aptinkamos mažiau
    nei puse atvejų.}
  \label{pav:prototipas}
\end{figure}

Prototipe išmatuota delsa didesnė nei \ref{tab:veikimas} lentelėje, nes
langai apdorojami po 250. Toks matavimas artimesnis diegimo sąlygoms ir
vis tiek lieka dviem eilėmis žemiau biudžeto.

